\section{Conclusion and future work}
\label{sec:conclusion}

We ported a predictive-maintenance EWC classifier onto a Cortex-M4 microcontroller and established, under a
paired protocol, an exact FP32 PC$\leftrightarrow$board parity in the frozen protocol, with an inference +
online learning latency $\ll$~\SI{100}{\milli\second} and a footprint below \SI{256}{\kilo\byte} (Gaps~2
and~3). On that same board we \emph{measured} the collapse of naive INT8 quantization of the binary head,
\emph{isolated its cause} through a bit-exact emulated ablation --- an uncalibrated per-tensor scale on too
narrow an accumulator --- then \emph{measured its recovery} with a calibrated kernel, with F1 back to
\num{0.9173} and \num{0.8995} and a frozen parity of \num{1.000} against the emulator. The diagnosis
remains emulated; the recovery no longer is.

Two complementary axes sharpen this result. On the quantization \emph{moment}, kernel calibration recovers
what matters and QAT training adds no decisive gain: when the training pipeline cannot be modified, a
properly calibrated PTQ suffices. On \emph{depth}, the boundary is not set by bits alone but by the
depth--metric pair: ternary holds in AUROC and breaks down in on-board measured F1, where the boundary
returns to INT4. Per-channel granularity pushes back the breakdown, and the memory gain advertised below
8 bits exists only with genuine
bit-packing --- which costs about \SI{55}{\micro\second} of unpacking, without threatening the budget.

\paragraph{What INT8 delivers, and what it does not.} The measured system budget decouples the three usual
promises. Weight RAM is indeed divided by~4, but total system RAM is unchanged (ratio \num{1.0002}),
because at this model scale the weights are not the dominant item; and that $\div 4$ holds for inference,
the floating head staying in memory as soon as the board learns. Latency \emph{increases} by about half
under inference and by a factor of~\num{2.4} under continual learning, requantization --- an item absent
from the floating-point path --- consuming one third of the inference budget on its own. Energy, finally, does not move: the measured gap is an order of magnitude below its uncertainty. On
this target INT8 is therefore justified by memory alone, and the energy lever lies elsewhere --- in the
clock frequency, where the latency margin we amply hold can be traded for battery life.

\paragraph{Future work.} The latency paradox points toward integer SIMD kernels (CMSIS-NN) for the MAC
item, and toward a power-of-two scale chain for the requantization item; this is the most direct route for
the memory gain to stop being paid in time. The learning regime calls for two routes of its own:
requantizing only the rows modified by the gradient step, which attacks the factor of~\num{2.4} directly,
and running gradient descent in the integer domain, the only way for the memory gain to bear on learning
itself rather than on inference alone. The "total RAM does not move" finding further calls for a
replication on a model where the weights \emph{do} dominate the footprint, a regime in which the $\div 4$
would become decisive again. Finally, three cells remain explicitly "to be measured": the \texttt{q15}
format on-board, the energy per online update, and the per-phase energy profile. We leave them empty
rather than estimate them --- the same rule that, throughout this article, separates what is measured
on-board from what is emulated on PC.
